\section{引言}\label{sec:intro}

\subsection{背景：受限算力下的代理替代}

F1 赛车的空气动力学设计高度依赖计算流体力学（CFD）与风洞实验。
2022 年 FIA 引入"地面效应"规则后，赛车设计面临两难：
一方面，海豚跳（Porpoising）等强非线性瞬态现象需要高保真仿真才能刻画；
另一方面，FIA 的空气动力学测试限制（ATR）对 CFD 时长与算力施加严格配额，
高保真仿真难以在有限周期内覆盖大规模参数空间。
在这一"高成本仿真 $+$ 受限算力"的约束下，使用\textbf{代理模型}
（在历史数据上训练的快速近似器）替代部分昂贵仿真并在其上做参数寻优，
成为工程上合理的路线。

\subsection{研究问题的定位}

本文不采用两种常见但在此并不成立的问题定义。

\vspace{0.2cm}
\noindent\textbf{方案 A：如何优化 F1 赛车的气动参数？}\quad
本研究使用的数据集并非真实 CFD 输出，而是由一个 XGBoost 预测引擎生成的数值近似
（详见第~\ref{sec:dataset}~节）。在此数据上得到的"最优参数"无法外推到真实赛车，
将研究定位为应用寻优会使结论缺乏物理可信度，故不采用。

\vspace{0.2cm}
\noindent\textbf{方案 B：如何提出一种更好的 PSO 算法？}\quad
本文引入了风险敏感粒子群优化，但它仅是整体框架中的一个防御模块；
且本文实验表明 PSO 并不显著优于 GA/DE（适应度极差仅 0.15），
若将主张落在"提出更好的优化器"上，将与本文发现不一致，故亦不采用。

\vspace{0.2cm}
\noindent
本文围绕如下研究问题展开：

\begin{chapterquestion}
\textbf{研究问题}：在受限计算预算下，如何构建一个可信（Trustworthy）的代理模型优化框架，
使其不仅给出最优解，也能评估该最优解的可信程度？
\end{chapterquestion}

研究关注点由"更快 / 更优 / 更高分"转向"可信性 / 可靠性 / 不确定性感知"。
当代理模型预测精度很高、优化结果却出现异常时，"准确"与"可信"之间的差异即为本文的研究对象。

\subsection{核心主张}

本文围绕可信代理优化问题展开，提出如下核心主张（Core Claim）：

\begin{chapterquestion}
\textbf{Core Claim}：在代理模型驱动的优化中，数据覆盖度（Data Coverage）、
不确定性管理（Uncertainty Management）与优化景观结构（Optimization Landscape），
比优化器选择本身更决定最终优化质量。
\end{chapterquestion}

围绕该主张，本文对各方法模块按其证据价值取舍：多模型对比仅保留选型依据，
SHAP 分析作为辅助解释模块。

\subsection{全文结构}

\begin{table}[H]
\centering
\caption{全文结构安排}
\begin{tabular}{ll}
\toprule
内容 & 对应章节 \\
\midrule
数据分析 & 第~\ref{sec:dataset}~节 \\
代理模型构建 & 第~\ref{sec:surrogate}~节 \\
异常现象观察 & 第~\ref{sec:anomaly}~节 \\
可信性分析（不确定性、风险敏感、数据精炼） & 第~\ref{sec:uncertainty}--\ref{sec:refine}~节 \\
景观诊断（解的趋同与综合分析） & 第~\ref{sec:converge}--\ref{sec:landscape}~节 \\
结论与讨论 & 第~\ref{sec:conclusion}~节 \\
\bottomrule
\end{tabular}
\end{table}

\subsection{研究意义}

在方法论层面，本文构建的"代理建模—不确定性感知—模型精炼—景观诊断"诊断流程，
可迁移至航空航天、船舶水动力学等具有"高成本仿真 $+$ 大规模参数空间 $+$ 数据不平衡"
特征的工程优化场景。其价值不在于给出某个具体的最优参数，
而在于为代理模型驱动的优化系统提供一套可信性诊断方法：
先评估数据与景观是否支撑优化，再讨论优化算法。

\subsection{本研究与计算智能课程主题的关系}

本研究的对象始终是\emph{代理模型驱动的智能优化}问题。
所采用的核心方法——粒子群优化（PSO）、遗传算法（GA）、差分进化（DE）
以及基于代理模型的优化（Surrogate Optimization）——均属计算智能的典型方法；
对这些算法的系统对比（见第~\ref{sec:converge}~节）构成本文的方法学基础。
不确定性估计与风险敏感优化并非偏离主题，而是用于提升计算智能优化系统在
分布外区域的\emph{可信性}。因此，本文是在计算智能框架内，
对"如何使智能优化结果更可信"这一问题的研究。
